\section{GIỚI THIỆU}
Ngày nay, các vấn đề về an toàn phòng cháy chữa cháy và rò rỉ khí độc hại tại các hộ gia đình, cơ sở sản xuất đang nhận được sự quan tâm rất lớn. Việc phát hiện sớm các nguy cơ này có vai trò cực kỳ quan trọng trong việc bảo vệ tính mạng và giảm thiểu thiệt hại tài sản. Hệ thống báo cháy tự động truyền thống thường chỉ có chức năng phát ra âm thanh báo động mà chưa có các cơ chế xử lý tức thời tự động như thông gió hay mở cửa thoát hiểm.

Nhằm giải quyết hạn chế trên, bài báo này đề xuất một giải pháp kết hợp giữa giám sát liên tục và phản hồi tự động dựa trên nền tảng vi điều khiển ESP8266 [1]. Bằng cách tích hợp nhiều cảm biến để đo nhiệt độ, nồng độ khí gas và phát hiện bức xạ từ ngọn lửa [2], hệ thống có khả năng đưa ra các quyết định điều khiển thiết bị ngoại vi ngay lập tức để can thiệp giảm thiểu rủi ro.

Phần còn lại của bài báo được tổ chức như sau: trong phần II, chúng tôi miêu tả chi tiết cấu trúc mô hình đề xuất. Trong phần III, chúng tôi đánh giá hiệu năng của hệ thống qua các kịch bản mô phỏng. Cuối cùng, chúng tôi kết luận bài báo trong phần IV.
